\documentclass[11pt]{article}

% Change "review" to "final" to generate the final (sometimes called camera-ready) version.
% Change to "preprint" to generate a non-anonymous version with page numbers.
% \usepackage[]{acl}
\usepackage[margin=1in]{geometry}
\usepackage{times}
\usepackage{latexsym}


% For proper rendering and hyphenation of words containing Latin characters (including in bib files)
\usepackage[T1]{fontenc}

% This assumes your files are encoded as UTF8
\usepackage[utf8]{inputenc}
\usepackage{microtype}
\usepackage{inconsolata}
\usepackage{graphicx}
\usepackage{url}
\usepackage{booktabs}
\usepackage{amsmath, amssymb}


\title{Group 42 Progress Step One:\\ Left and Right Leaning Bias Detection in Media Headlines}


\author{
  Ahren Chen, Fei Xie, Grace Xiao\\
  \texttt{\{chena125, xief17, xiaot13\}@mcmaster.ca}
}
\begin{document}
\maketitle

\section{Annotation Instructions}

\section{Dataset}
\subsection{Description}

\subsection{Interesting Data Points}

\subsection{Annotation Agreement}

\section{Discussion Questions}
\begin{enumerate}
  \item What did you learn about the task from doing the annotation?

  \item What challenges do you expect models to face when learning from your data?
  \item What surprising things did you observe in the data?
  \item Which features do you expect to be useful
  \item Describe the mental model you came up with to do the task at a high level. What
process did you use?
  \item Was there anything unclear about the instructions? How would you recommend modi-
fying them in the future to make them more helpful?
\end{enumerate}


\end{document}
